\documentclass[12pt]{article}
\usepackage[utf8]{inputenc}
\usepackage{amsmath, amssymb, geometry}
\geometry{letterpaper, margin=1in}

\title{\textbf{Steps to Calculate the Launch Window using a Pork Chop Plot}}
\author{}
\date{}

\begin{document}

\maketitle

This detailed approach moves systematically from defining parameters to analyzing the final visualization, breaking down the complex computational process into manageable stages, conceptually mimicking a professional brainstorming session.

\begin{enumerate}
    \item \textbf{Define Mission Parameters \& Constraints:}
    \begin{itemize}
        \item Identify the launch planet (e.g., Earth).
        \item Identify the target planet (e.g., Mars).
        \item Establish the overall search window dates, defining ranges for potential departure dates ($t_{dep}$) and potential arrival dates ($t_{arr}$). Consider realistic mission timelines and planetary alignments.
        \item Define relevant mission constraints (e.g., maximum launch vehicle performance ($C_3$ limits), acceptable Time of Flight (TOF) ranges, arrival velocity constraints, etc.).
    \end{itemize}

    \item \textbf{Initialize the Grid Search:}
    \begin{itemize}
        \item Create a detailed 2D grid of potential departure dates and potential arrival dates within the defined search windows. This grid will populate the entire porkchop plot area.
        \item Determine an appropriate step size (e.g., every 2 days) for both axes to step through sequentially, balancing resolution and computational time.
        \item \textit{Conceptual Brainstorming Note:} This is where we create our map structure before exploring!
    \end{itemize}

    \item \textbf{Perform the Lambert Sweep (Iterative Loop):} \\
    For every single combination of ($t_{dep}$, $t_{arr}$) in your initialized grid, the computer will iteratively perform the following sequential calculations:
    \begin{enumerate}
        \item \textbf{Calculate Precise Positions:} Use high-accuracy planetary ephemerides (e.g., NASA SPICE data kernels) to determine the exact heliocentric position vectors of the launch planet ($\mathbf{r}_1$ at $t_{dep}$) and the target planet ($\mathbf{r}_2$ at $t_{arr}$).
        \item \textbf{Determine Time of Flight (TOF):} Calculate the Time of Flight, $TOF = t_{arr} - t_{dep}$. (Ensure $TOF$ is positive; impossible/negative flight times are conceptually discarded or handled by the loop logic).
        \item \textbf{Solve Lambert's Boundary Value Problem:} Execute a Lambert solver numerical algorithm (like the Izzo solver) using the Sun's gravitational parameter, $\mathbf{r}_1$, $\mathbf{r}_2$, and $TOF$. This mathematically calculates the precise elliptical transfer orbit arc connecting the two planets in that specific time, outputting the necessary heliocentric velocity vectors ($\mathbf{v}_1$ at departure and $\mathbf{v}_2$ at arrival). For initial search, typically focus on direct, 0-revolution transfers.
        \item \textbf{Calculate Hyperbolic Excess Velocity ($\mathbf{v}_{\infty}$):} Compute the $\mathbf{v}_{\infty}$ vector at departure, which is the vector difference between the heliocentric transfer velocity ($\mathbf{v}_1$) and the launch planet's actual heliocentric velocity at $t_{dep}$ ($\mathbf{v}_{planet\_1}$).
        \item \textbf{Calculate $C_3$ Energy:} Compute the Characteristic Energy ($C_3$) by squaring the magnitude of the Hyperbolic Excess Velocity vector: $C_3 = | \mathbf{v}_{\infty} |^2$. (Conceptually, this is the energy cost associated with leaving the planet!)
        \item \textbf{Store Calculated Data:} Systematically store the calculated $C_3$ value (and optionally other parameters like $\Delta V$, arrival $v_{\infty}$, DLA, TOF, etc.) associated specifically with the current ($t_{dep}$, $t_{arr}$) grid point.
    \end{enumerate}

    \item \textbf{Generate the Pork Chop Plot Visualization:}
    \begin{itemize}
        \item Once the computational sweep is complete for the entire grid, use data visualization software (like Matplotlib) to plot the stored $C_3$ values as a contour map.
        \item \textit{Brainstorming Layout:} X-axis represents 'Departure Date', Y-axis represents 'Arrival Date', and contour lines indicate $C_3$ levels.
        \item Apply dynamic filtering and shading – low-energy regions (cool colors/blue) are clearly highlighted, while high-energy/unfeasible regions are either filtered out or rendered in warm colors.
    \end{itemize}

    \item \textbf{Analyze \& Identify the Window of Opportunity:}
    \begin{itemize}
        \item Visually inspect the generated Pork Chop plot. The characteristic "porkchop" shapes with distinct low-energy centers (valleys) will become evident.
        \item Identify these dark blue/cool-colored centers – they represent the multiple potential optimal launch opportunity windows.
        \item Select a specific trajectory date combination within one of these valleys, conceptually trading off lowest $C_3$ energy (most payload) against desirable Time of Flight, feasible launch vehicle performance limits, and whether a Type I ($\Delta \nu < 180^\circ$) or Type II ($\Delta \nu > 180^\circ$) trajectory is preferred for mission objectives. \textit{Example:} For Mars 2020, designers used Porkchop plots to select a lower energy Type II window to maximize payload!
    \end{itemize}
\end{enumerate}

\end{document}